\chapter{Polynomials}\label{ch:poly}

Collecting like terms is the first algebra anyone learns and the first rule that needed a real
proof: $3x + 5x = 8x$ is the distributive law, and the distributive law is an axiom of a ring.
Expansion --- multiplying out $(x + y)^{12}$ --- is the same law applied many times, and the
chapter's second half is about the fact that ``many times'' can be organized so that the answer
is a polynomial rather than a mess.

\section{Rings and the distributive law}

\begin{definition}[Commutative ring]
  A commutative ring is a set with two operations $+$ and $\cdot$ such that $+$ is associative
  and commutative with an identity $0$ and inverses, $\cdot$ is associative and commutative with
  an identity $1$, and multiplication distributes over addition:
  $a(b + c) = ab + ac$.
\end{definition}

Everything a simplifier does to a sum or product follows from these axioms. The identities
$a + 0 = a$ and $a \cdot 1 = a$ are the rule \rn{simp.identity}; $a \cdot 0 = 0$ is a
\emph{theorem} of rings, not an axiom, and its proof is worth remembering:
$a \cdot 0 = a \cdot (0 + 0) = a \cdot 0 + a \cdot 0$, and cancelling $a \cdot 0$ from both sides
gives $a \cdot 0 = 0$. Associativity of $+$ and $\cdot$ is \rn{simp.flatten}, which turns a nested
sum into one flat list.

Like terms are the distributive law read right to left:
\[
  a t + b t = (a + b)\, t .
\]
Here $t$ is any term and $a, b$ are numerals. The rule is stated for two terms; a sum of $k$ like
terms collects by $k - 1$ applications.

\section{The polynomial ring}

\begin{definition}[Polynomials]
  For a commutative ring $R$, the polynomial ring $R[x_1, \ldots, x_n]$ is the set of finite
  formal sums $\sum_\alpha c_\alpha x^\alpha$, where $\alpha$ ranges over exponent vectors in
  $\mathbb{N}^n$, $x^\alpha = x_1^{\alpha_1} \cdots x_n^{\alpha_n}$, and each $c_\alpha \in R$.
  Addition is coefficient-wise; multiplication is
  \[
    \Big(\sum_\alpha a_\alpha x^\alpha\Big)\Big(\sum_\beta b_\beta x^\beta\Big)
      = \sum_\gamma \Big(\sum_{\alpha + \beta = \gamma} a_\alpha b_\beta\Big) x^\gamma .
  \]
\end{definition}

Two facts about this definition matter for the engine. First, a polynomial is a \emph{function
from monomials to coefficients} with finite support; the monomials $x^\alpha$ are a basis, and two
polynomials are equal iff their coefficients agree monomial by monomial. That is what makes
``collect like terms'' well defined: the like terms are the summands with the same monomial.
Second, multiplication is the distributive law applied to every pair of summands, followed by
collection. So expansion is not a search; it is a computation with one answer.

\begin{theorem}[Binomial theorem]
  For $n \in \mathbb{N}$, $(x + y)^n = \sum_{k=0}^n \binom{n}{k} x^k y^{n-k}$.
\end{theorem}
\begin{proof}
  Induction on $n$. For $n = 0$ both sides are $1$. Assuming the identity for $n$,
  \[
    (x+y)^{n+1} = (x+y)\sum_{k} \binom{n}{k} x^k y^{n-k}
      = \sum_k \binom{n}{k} x^{k+1} y^{n-k} + \sum_k \binom{n}{k} x^k y^{n+1-k},
  \]
  and collecting the coefficient of $x^k y^{n+1-k}$ gives $\binom{n}{k-1} + \binom{n}{k} =
  \binom{n+1}{k}$, Pascal's rule.
\end{proof}

The proof is itself an expansion followed by a collection, which is the engine's algorithm with
the coefficients tracked symbolically. Pascal's rule is the collection step.

\section{The rule: collecting like terms}

A summand is read as a coefficient times a rest, \lc{coeffRest}: $3x$ is $(3, x)$, $x$ is
$(1, x)$, $5$ is $(5, 1)$. The rule finds the first pair of summands with equal rests and merges
them:

\begin{minted}{lean}
def mergeTerms : List Expr → Option (List Expr × Expr)
  | [] => none
  | e :: rest =>
    let (c, t) := coeffRest e
    if bigBase t then
      match rest.find? (fun f => equal (coeffRest f).2 t) with
      | some f => some ((.mul [.num (c + (coeffRest f).1), t]) :: removeFirst (fun f => equal (coeffRest f).2 t) rest, t)
      | none => (mergeTerms rest).map fun (l, u) => (e :: l, u)
    else (mergeTerms rest).map fun (l, u) => (e :: l, u)
\end{minted}

The guard \lc{bigBase t} excludes a rest that is a bare numeral, so that $2 + 3$ is left to the
fold-constants rule of Chapter~\ref{ch:q} rather than to this one; the two rules divide the work
by the shape of the term, and the termination proof needs the division to be clean.

\begin{thmbox}[\lcn{collectTerms_soundR}]
  Wherever the rule fires, \lc{evalR} of the output equals \lc{evalR} of the input.
\end{thmbox}

The proof is the distributive law of $\mathbb{R}$ once (\lc{mergeTerms_soundR}: the merged
summand has the value of the two it replaced) and the permutation lemma for sums once (the
rule moved the merged summand to the front). Nothing else. The status is \status{verified}.

\section{Expansion as a function}

Expansion could have been a rule --- ``a product with a sum among its factors rewrites to the
sum of the products'' --- and in the first version of the engine it was. Chapter~\ref{ch:cannot}
explains why it cannot coexist with $x \cdot x \to x^2$ under one termination ordering. Here is
the mathematics that replaced it: the polynomial ring's multiplication, computed directly.

\begin{minted}{lean}
structure Mono where
  coeff : Q
  key : List Expr     -- the non-numeral factors, sorted

abbrev Poly := List Mono   -- distinct keys

def insertMono (m : Mono) : Poly → Poly
  | [] => [m]
  | n :: ns => if Expr.beqList n.key m.key then ⟨n.coeff + m.coeff, n.key⟩ :: ns else n :: insertMono m ns

def polyMul (p q : Poly) : Poly := collect (p.flatMap fun a => q.map (Mono.mul a))

def distMul (fs : List Expr) : Expr :=
  if fs.any isAdd then ofPoly ((fs.map toPoly).foldl polyMul [⟨Q.one, []⟩]) else .mul fs
\end{minted}

Read against the definition: a \lc{Mono} is one term $c_\alpha x^\alpha$, with the monomial
represented by its sorted factor list so that equal monomials have equal keys; a \lc{Poly} is
the finite-support function, as a list with distinct keys; \lc{insertMono} is coefficient-wise
addition; \lc{polyMul} is the double sum in the definition, with the inner collection done by
\lc{collect}. The function \lc{dist} applies this bottom-up, and a power of a sum with a small
natural exponent becomes that many copies multiplied out.

The collection \emph{during} multiplication is what makes the binomial theorem's $n = 12$
tractable. Without it, $(x+y)^{12}$ produces $2^{12} = 4096$ monomials and leaves them for the
collect-like-terms rule, which merges one pair at a time; with it, the $13$ monomials of the
answer are the only ones that ever exist.

\section{The theorem for expansion}

The soundness proof is a structural induction on \lc{dist}, and it is stated for any semantics
that reads sums, products, numerals and natural powers the obvious way:

\begin{minted}{lean}
structure Sem where
  ev : EnvR → Expr → ℝ
  ev_add : ∀ ρ es, ev ρ (.add es) = (es.map (ev ρ)).sum
  ev_mul : ∀ ρ es, ev ρ (.mul es) = (es.map (ev ρ)).prod
  ev_num : ∀ ρ q, ev ρ (.num q) = (q.val : ℝ)
  ev_pow : ∀ ρ b e, ev ρ (.pow b e) = ev ρ b ^ ev ρ e

def evalM (m : Mono) : ℝ := (m.coeff.val : ℝ) * (m.key.map (S.ev ρ)).prod
def evalP : Poly → ℝ
  | [] => 0
  | m :: ms => evalM S ρ m + evalP ms

theorem dist_sound (ρ : EnvR) : ∀ e : Expr, evalR ρ (Expand.dist e) = evalR ρ e
\end{minted}

The chain of lemmas is the chain of definitions: a monomial evaluates to its coefficient times
its factors (\lc{ev_Mono_eval}); inserting a monomial adds its value (\lc{evalP_insertMono},
which is the distributive law $c_1 t + c_2 t = (c_1 + c_2) t$ once more); the product of two
polynomials evaluates to the product of their values (\lc{evalP_polyMul}, distributivity over
the double sum); and \lc{dist} is sound by induction. The structure \lc{Sem} is instantiated
twice, for \lc{evalR} here and for the derivative semantics \lc{evalD} in
Chapter~\ref{ch:parts}, where expansion is part of the integral checker. The status of
\rn{cmd.expand} is \status{verified}, unconditionally: distribution holds in every commutative
ring, and $\mathbb{R}$ is one.

\begin{trybox}
\begin{minted}{text}
3*x + 5*x + x
expand((x + y)^12)
expand((a + b)*(a - b))
\end{minted}
  The coefficient of $x^6 y^6$ in the second is $\binom{12}{6} = 924$. The third is
  $a^2 - b^2$ after the pipeline collects $ab - ba$.
\end{trybox}

\section{Exercises}

\begin{exercisebox}[Pascal's rule]
  Prove $\binom{n}{k-1} + \binom{n}{k} = \binom{n+1}{k}$ from the definition
  $\binom{n}{k} = \frac{n!}{k!(n-k)!}$, and then combinatorially.
\end{exercisebox}

\begin{exercisebox}[Keys]
  Why must the key of a monomial be \emph{sorted}? Give two factor lists for the same monomial
  that \lc{Expr.beqList} would distinguish if they were not.
\end{exercisebox}

\begin{exercisebox}[Multinomial]
  State and prove the multinomial theorem for $(x + y + z)^n$ by induction on $n$ using the
  binomial theorem, and count the monomials of $(x + y + z)^{4}$.
\end{exercisebox}

\begin{exercisebox}[The generic semantics]
  \lc{Sem} does not mention functions or variables. Explain why the soundness of \lc{dist} does
  not need to know what $\sin$ means.
\end{exercisebox}
